\Askhsh{34320}{4}
\begin{ylh}{1}
    \item 3.1. Είδη και στοιχεία τριγώνων
    \item 4.2. Τέμνουσα δύο ευθειών - Ευκλείδειο αίτημα
    \item 4.6. Άθροισμα γωνιών τριγώνου
    \item 5.10. Τραπέζιο.
\end{ylh}
% ===================================================================
%  Αρχείο: 34320.tex (Fragment)
%  Αυτόματη μετατροπή μέσω doc2latex.py
% ===================================================================

ΘΕΜΑ 4

Έστω Α, Β και Γ συνευθειακά σημεία με ΑΒ = 2ΒΓ. Θεωρούμε το μέσο Μ της ΑΒ. Προς το ίδιο ημιεπίεδο κατασκευάζουμε τα ισόπλευρα τρίγωνα ΑΔΒ και ΒΕΓ.

Να αποδείξετε ότι:

\begin{alist}
\item το τετράπλευρο ΑΔΕΒ είναι τραπέζιο με βάσεις τα τμήματα ΑΔ και ΒΕ, \monades{9}

\item τα τρίγωνα ΔΜΒ και ΔΕΒ είναι ίσα, \monades{8}

\item Δ$\widehat{Μ}$Β + Δ$\widehat{Ε}$Β = 180\textsuperscript{0}. \monades{8}


\begin{center}
\begin{tikzpicture}[scale=0.8]
\tkzDefPoint(0,0){A}
\tkzDefPoint(4,0){B}
\tkzDefPoint(6,0){C}
\tkzDefMidPoint(A,B) \tkzGetPoint{M}

\tkzDefTriangle[equilateral](A,B) \tkzGetPoint{D}
\tkzDefTriangle[equilateral](B,C) \tkzGetPoint{E}

\draw[l3] (A)--(C);
\draw[l3] (A)--(D)--(B)--cycle;
\draw[l3] (B)--(E)--(C)--cycle;
\draw[l3, dashed] (D)--(E);
\draw[l3, dashed] (D)--(M);

\tkzDrawPoints(A,B,C,M,D,E)
\tkzLabelPoint[below](A){$A$}
\tkzLabelPoint[below](B){$B$}
\tkzLabelPoint[below](C){$\Gamma$}
\tkzLabelPoint[below](M){$M$}
\tkzLabelPoint[above](D){$\Delta$}
\tkzLabelPoint[above right](E){$E$}

\tkzMarkSegments[mark=s|](A,M M,B B,C B,E E,C)
\tkzMarkSegments[mark=s||](A,D D,B)
\end{tikzpicture}
\end{center}
\end{alist}
